\documentclass[11pt,a4paper]{article}

\usepackage[margin=2.5cm]{geometry}
\usepackage[T1]{fontenc}
\usepackage[utf8]{inputenc}
\usepackage{lmodern}
\usepackage{microtype}
\usepackage{graphicx}
\usepackage{float}
\usepackage{xcolor}
\usepackage{listings}
\usepackage{hyperref}

\hypersetup{
  colorlinks=true,
  linkcolor=black,
  urlcolor=blue,
  pdfauthor={Oriol Freixa},
  pdftitle={Essentia Danceability Plugin Report}
}

\lstset{
  basicstyle=\ttfamily\small,
  breaklines=true,
  frame=single,
  columns=fullflexible,
  backgroundcolor=\color{black!3}
}

\title{Essentia Danceability Plugin}
\author{Oriol Freixa}
\date{\today}

\begin{document}

\maketitle

\section{Overview}

We followed the Essentia plugins setup, built the required JUCE/Essentia
dependencies, and created a new Projucer plugin from the RMS template. We
renamed the project to \texttt{Danceability} and adapted the RMS analyzer into
a VST3 analyzer based on Essentia's \texttt{Danceability} algorithm.

The source code is available in the GitHub repository
\href{https://github.com/OriolFreixa/essentia_plugin_actsm}{OriolFreixa/essentia\_plugin\_actsm}.

\begin{figure}[H]
  \centering
  \includegraphics[width=0.82\linewidth]{\detokenize{../imgs_and_vids/Screenshot 2026-06-14 at 20.06.04.png}}
  \caption{Danceability plugin loaded in Ableton Live.}
\end{figure}

\section{Implementation}

We replaced the RMS algorithm with \texttt{Danceability}. Since this descriptor
needs a longer signal than a normal DAW audio callback, we adapted the
processor so that it stored a rolling mono analysis window and computed the
descriptor periodically instead of per block.

\begin{lstlisting}[language=C++]
danceability = AlgorithmFactory::create("Danceability");
danceability->input("signal").set(essentiaBuffer);
danceability->output("danceability").set(danceabilityValue);
danceability->output("dfa").set(dfaValues);
\end{lstlisting}

We connected both Essentia outputs: the scalar \texttt{danceability} value is
shown in the UI, while the \texttt{dfa} vector is kept internally because it is
part of the algorithm output but not useful as a main user-facing value.

\begin{lstlisting}[language=C++]
essentiaBuffer.resize(previousSize + numSamples);
// Mix host input to mono and append it to the rolling window.

if (essentiaBuffer.size() > analysisWindowSamples)
    essentiaBuffer.erase(essentiaBuffer.begin(),
                         essentiaBuffer.begin() + samplesToDrop);

if (essentiaBuffer.size() >= analysisWindowSamples)
    danceability->compute();
\end{lstlisting}

\section{Parameters}

We exposed \texttt{minTau}, \texttt{maxTau}, and \texttt{tauMultiplier} because
they are real Essentia Danceability parameters and make the analysis tunable
from the host. We used \texttt{AudioProcessorValueTreeState} so the values can
be automated and saved with the session.

\begin{lstlisting}[language=C++]
params.push_back(std::make_unique<juce::AudioParameterFloat>(
    juce::ParameterID {"minTau", 1}, "minTau",
    juce::NormalisableRange<float> {50.f, 1000.f, 10.f}, 310.f));

params.push_back(std::make_unique<juce::AudioParameterFloat>(
    juce::ParameterID {"maxTau", 1}, "maxTau",
    juce::NormalisableRange<float> {1000.f, 8800.f, 100.f}, 8800.f));

params.push_back(std::make_unique<juce::AudioParameterFloat>(
    juce::ParameterID {"tauMultiplier", 1}, "tauMultiplier",
    juce::NormalisableRange<float> {1.01f, 1.5f, 0.01f}, 1.1f));
\end{lstlisting}

We also added \texttt{Analysis Window} as a useful extra parameter. It is not
an Essentia Danceability parameter, but it controls how many seconds of recent
audio we keep before analysis, which makes the plugin more practical in a DAW.

\begin{lstlisting}[language=C++]
params.push_back(std::make_unique<juce::AudioParameterFloat>(
    juce::ParameterID {"analysisWindowSeconds", 1}, "Analysis Window",
    juce::NormalisableRange<float> {1.f, 15.f, 0.1f}, 10.f));

const auto analysisWindowSeconds =
    apvts.getRawParameterValue("analysisWindowSeconds")->load();

analysisWindowSamples =
    static_cast<size_t>(currentSampleRate * analysisWindowSeconds);
\end{lstlisting}

\section{Interface And Test}

We replaced the RMS labels with a compact Danceability Analyzer interface,
added sliders for the exposed parameters, and kept the UPF and Essentia logos.
We tested the plugin in Ableton Live using the included test project and
recorded a short demo video. The video is available in
\texttt{imgs\_and\_vids/danceability-demo.mov}.

\begin{figure}[H]
  \centering
  \includegraphics[width=0.82\linewidth]{\detokenize{../imgs_and_vids/Screenshot 2026-06-14 at 20.18.24.png}}
  \caption{Analyzer running and showing a non-zero Danceability value.}
\end{figure}

\end{document}
